\documentclass[preview]{standalone}
\usepackage[english]{babel}
\usepackage{amsmath}
\usepackage{amssymb}
\usepackage{ragged2e}
\setlength\parindent{0pt}
\begin{document}
\begin{center}
\special{dvisvgm:raw <g id='unique000'>}\parbox{9cm}{\noindent{}\justifying{\special{dvisvgm:raw </g>}\special{dvisvgm:raw <g id='unique001'>}\begin{itemize} \item A group is unimodular if the modular function has a value of 1 for every element.\vspace{-10pt}\end{itemize}\special{dvisvgm:raw </g>}\special{dvisvgm:raw <g id='unique002'>}\begin{itemize} \item In a totally disconnected locally compact group we have that $\Delta(g) = \frac{s(g)}{s(g^{-1})}$\vspace{-10pt}.\end{itemize}\special{dvisvgm:raw </g>}\special{dvisvgm:raw <g id='unique003'>}\begin{itemize} \item For a group with compact-open vertex stabilisers, a unimodularity condition was derived by Bass-Kulkarni in [BK90]\vspace{-10pt}.\end{itemize}\special{dvisvgm:raw </g>}\special{dvisvgm:raw <g id='unique004'>}\begin{itemize} \item For (P)-closed groups, the same condition can be found by using our formula for the scale function.\vspace{-10pt}\end{itemize}\special{dvisvgm:raw </g>}\special{dvisvgm:raw <g id='unique005'>}}}\special{dvisvgm:raw </g>}
\end{center}
\end{document}